\chapter{Table of Identities}
\label{app:identities}

The identities developed across Chapters~8--10 are collected here for
reference, grouped by family. Each entry names the section in which the
identity was derived, since every identity in this book is a consequence
of prior results rather than an independent fact to memorize.

\section{Reciprocal and Quotient Identities}

\begin{align*}
\csc\theta &= \frac{1}{\sin\theta} & \sec\theta &= \frac{1}{\cos\theta} &
\cot\theta &= \frac{1}{\tan\theta} \\
\tan\theta &= \frac{\sin\theta}{\cos\theta} & \cot\theta &=
\frac{\cos\theta}{\sin\theta} & & \text{(Chapter 9)}
\end{align*}

\section{Pythagorean Identities}

\begin{align*}
\sin^2\theta + \cos^2\theta &= 1 & &\text{(Chapter 9)} \\
1 + \tan^2\theta &= \sec^2\theta & &\text{(Chapter 9)} \\
1 + \cot^2\theta &= \csc^2\theta & &\text{(Chapter 9)}
\end{align*}

\section{Cofunction Identities}

\[
\sin\left(\frac{\pi}{2} - \theta\right) = \cos\theta, \qquad
\cos\left(\frac{\pi}{2} - \theta\right) = \sin\theta, \qquad
\tan\left(\frac{\pi}{2} - \theta\right) = \cot\theta. \qquad \text{(Chapter
9)}
\]

\section{Even and Odd Identities}

\[
\cos(-\theta) = \cos\theta, \qquad \sin(-\theta) = -\sin\theta, \qquad
\tan(-\theta) = -\tan\theta. \qquad \text{(Chapters 6, 9)}
\]

\section{Sum and Difference Formulas}

\begin{align*}
\cos(A+B) &= \cos A\cos B - \sin A\sin B & \sin(A+B) &= \sin A\cos B +
\cos A\sin B \\
\cos(A-B) &= \cos A\cos B + \sin A\sin B & \sin(A-B) &= \sin A\cos B -
\cos A\sin B
\end{align*}
\begin{center}(Chapter 8)\end{center}

\section{Double-Angle Formulas}

\[
\sin(2\theta) = 2\sin\theta\cos\theta, \qquad
\cos(2\theta) = \cos^2\theta - \sin^2\theta = 2\cos^2\theta - 1 =
1 - 2\sin^2\theta, \qquad
\tan(2\theta) = \frac{2\tan\theta}{1-\tan^2\theta}. \qquad \text{(Chapter
10)}
\]

\section{Half-Angle Formulas}

\[
\sin\left(\frac{\theta}{2}\right) = \pm\sqrt{\frac{1-\cos\theta}{2}},
\qquad
\cos\left(\frac{\theta}{2}\right) = \pm\sqrt{\frac{1+\cos\theta}{2}},
\qquad
\tan\left(\frac{\theta}{2}\right) =
\pm\sqrt{\frac{1-\cos\theta}{1+\cos\theta}}. \qquad \text{(Chapter 10)}
\]
The sign in each formula is determined by the quadrant of $\theta/2$, not
the quadrant of $\theta$; see Section~10.4.

\section{Product-to-Sum Formulas}

\begin{align*}
\sin A\cos B &= \frac{\sin(A+B)+\sin(A-B)}{2} \\
\cos A\cos B &= \frac{\cos(A+B)+\cos(A-B)}{2} \\
\sin A\sin B &= \frac{\cos(A-B)-\cos(A+B)}{2}
\end{align*}
\begin{center}(Chapter 10)\end{center}

\section{Sum-to-Product Formulas}

\[
\sin x + \sin y = 2\sin\left(\frac{x+y}{2}\right)\cos\left(\frac{x-y}{2}\right),
\qquad
\cos x + \cos y = 2\cos\left(\frac{x+y}{2}\right)\cos\left(\frac{x-y}{2}\right).
\qquad \text{(Chapter 10)}
\]

\section{Law of Sines and Law of Cosines}

\[
\frac{a}{\sin A} = \frac{b}{\sin B} = \frac{c}{\sin C}, \qquad
c^2 = a^2+b^2-2ab\cos C. \qquad \text{(Chapter 12)}
\]

\section{Exact Values of Common Angles}

\begin{table}[h]
\centering
\begin{tabular}{ccccc}
\toprule
$\theta$ (degrees) & $0^\circ$ & $30^\circ$ & $45^\circ$ & $60^\circ$
\quad $90^\circ$ \\
\midrule
$\theta$ (radians) & $0$ & $\pi/6$ & $\pi/4$ & $\pi/3$ \quad $\pi/2$ \\
$\sin\theta$ & $0$ & $1/2$ & $\sqrt2/2$ & $\sqrt3/2$ \quad $1$ \\
$\cos\theta$ & $1$ & $\sqrt3/2$ & $\sqrt2/2$ & $1/2$ \quad $0$ \\
$\tan\theta$ & $0$ & $\sqrt3/3$ & $1$ & $\sqrt3$ \quad undefined \\
\bottomrule
\end{tabular}
\caption{Exact trigonometric values derived in Chapter~3.}
\end{table}

\section{Euler's Formula and Related Identities}

\[
e^{i\theta} = \cos\theta + i\sin\theta, \qquad
\cos\theta = \frac{e^{i\theta}+e^{-i\theta}}{2}, \qquad
\sin\theta = \frac{e^{i\theta}-e^{-i\theta}}{2i}. \qquad \text{(Chapter
16)}
\]
De Moivre's theorem: $(\cos\theta+i\sin\theta)^n = \cos(n\theta) +
i\sin(n\theta)$. \quad (Chapter 16)
